% ============================================================
%  J6 — IA, Éthique & Société
%  Julien Rolland — M2 Développement Fullstack
% ============================================================
\documentclass[aspectratio=169, 10pt]{beamer}
\input{../shared/preamble}

\title[Éthique \& IA]{IA, Éthique \& Société}
\subtitle{Jour 6}
\date{Jour 6}

% ============================================================
\begin{document}
% ============================================================

\begin{frame}
  \titlepage
\end{frame}

\begin{frame}{Plan du module}
  \tableofcontents
\end{frame}

% ============================================================
\section{Pourquoi l'éthique en IA maintenant ?}
% ============================================================

\begin{frame}{L'IA prend des décisions réelles}
  \begin{columns}[T]
    \begin{column}{0.52\textwidth}
      \begin{alertblock}{Ce n'est plus de la recherche}
        Des systèmes d'IA décident \textbf{aujourd'hui} :
        \begin{itemize}
          \item Si votre crédit est accordé
          \item Si votre CV est lu par un recruteur
          \item Quelle peine de prison recommander
          \item Quelle publicité vous montrer
          \item Si une tumeur est maligne
        \end{itemize}
      \end{alertblock}
    \end{column}
    \begin{column}{0.44\textwidth}
      \begin{block}{Une asymétrie de pouvoir}
        \begin{itemize}
          \item La personne affectée ne sait souvent
                \textbf{pas qu'un algorithme} a décidé
          \item Elle ne peut ni \textbf{contester}
                ni \textbf{comprendre} la décision
          \item Les erreurs se \textbf{reproduisent à l'échelle}
                — pas une personne, des millions
        \end{itemize}
      \end{block}
      \smallskip
      \begin{exampleblock}{La question centrale}
        Qui est responsable quand un algorithme
        se trompe — et que cette erreur coûte cher
        à quelqu'un ?
      \end{exampleblock}
    \end{column}
  \end{columns}
\end{frame}

\begin{frame}{Trois cas qui ont marqué les esprits}
  \begin{columns}[T]
    \begin{column}{0.31\textwidth}
      \begin{alertblock}{COMPAS --- Justice pénale (2016)}
        Algorithme prédit le risque de récidive
        aux États-Unis.\\[4pt]
        Étude ProPublica : les faux positifs
        («~va récidiver~» alors que non)
        étaient \textbf{deux fois plus fréquents}
        pour les prévenus noirs que blancs.
      \end{alertblock}
    \end{column}
    \begin{column}{0.31\textwidth}
      \begin{alertblock}{Amazon Hiring (2018)}
        Outil de tri de CV entraîné sur
        les embauches des 10 années précédentes.\\[4pt]
        Résultat : pénalisait automatiquement
        les CV contenant le mot \textbf{«~femmes~»}
        (ex: «~capitaine de l'équipe féminine~»).\\[4pt]
        Arrêté en interne avant déploiement.
      \end{alertblock}
    \end{column}
    \begin{column}{0.31\textwidth}
      \begin{alertblock}{Google Translate --- Genre (2017)}
        Traduction du turc (sans genre grammatical)
        vers l'anglais.\\[4pt]
        «~O bir doktor~» (il/elle est médecin)
        $\to$ \textbf{«~He is a doctor~»}\\[2pt]
        «~O bir hemşire~» (il/elle est infirmier/e)
        $\to$ \textbf{«~She is a nurse~»}\\[4pt]
        Le modèle reproduit les stéréotypes
        présents dans le corpus d'entraînement.
      \end{alertblock}
    \end{column}
  \end{columns}
\end{frame}

% ============================================================
\section{Les biais}
% ============================================================

\begin{frame}{D'où vient un biais ?}
  \begin{center}
    \begin{tikzpicture}[font=\small,
      blk/.style={draw, rounded corners=4pt, align=center,
                  minimum width=2.2cm, minimum height=0.9cm,
                  fill=jedy_blue!10, draw=jedy_blue},
      bad/.style={draw, rounded corners=4pt, align=center,
                  minimum width=2.0cm, minimum height=0.7cm,
                  fill=jedy_alert!20, draw=jedy_alert},
      arr/.style={->, thick, jedy_mid}]
      \node[blk] (data)  at (0,   0) {Données};
      \node[blk] (model) at (4,   0) {Modèle};
      \node[blk] (pred)  at (8,   0) {Décision};
      \node[blk] (world) at (-3,  0) {Monde réel};
      \draw[arr] (world) -- (data);
      \draw[arr] (data)  -- (model);
      \draw[arr] (model) -- (pred);
      \node[bad] at (0,  -1.5) {Biais de\\représentation};
      \node[bad] at (0,  -2.5) {Biais historique};
      \node[bad] at (4,  -1.5) {Biais de proxy};
      \node[bad] at (4,  -2.5) {Choix de la loss};
      \node[bad] at (8,  -1.5) {Biais\\d'interprétation};
      \draw[->, dashed, jedy_alert] (0, -1.2) -- (data.south);
      \draw[->, dashed, jedy_alert] (4, -1.2) -- (model.south);
      \draw[->, dashed, jedy_alert] (8, -1.2) -- (pred.south);
    \end{tikzpicture}
  \end{center}
  \vspace{4pt}
  \begin{block}{Idée clé}
    Un biais peut s'introduire \textbf{à chaque étape} du pipeline —
    même avec les meilleures intentions.
    Avoir un bon score de test ne garantit pas l'absence de biais.
  \end{block}
\end{frame}

\begin{frame}{Biais de représentation}
  \begin{columns}[T]
    \begin{column}{0.52\textwidth}
      \begin{alertblock}{Quand les données ne reflètent pas le monde}
        Si certaines populations sont \textbf{sous-représentées}
        dans le dataset, le modèle apprend moins bien pour elles —
        même si la performance globale semble bonne.\\[4pt]
        Exemples :
        \begin{itemize}
          \item Reconnaissance faciale entraînée sur
                des visages majoritairement blancs
          \item Textes sur internet sur-représentant
                certains groupes, langues, cultures
        \end{itemize}
      \end{alertblock}
    \end{column}
    \begin{column}{0.44\textwidth}
      \begin{block}{Mécanisme}
        Le modèle optimise pour minimiser l'erreur
        sur l'ensemble du dataset.\\[4pt]
        Les groupes majoritaires «~pèsent~» davantage
        dans la loss — le modèle leur donne
        implicitement la priorité.
      \end{block}
      \smallskip
      \begin{exampleblock}{Conséquence}
        Un modèle avec 95\,\% d'accuracy globale
        peut n'avoir que 70\,\% sur un sous-groupe
        minoritaire — sans que ça se voie
        dans la métrique globale.
      \end{exampleblock}
    \end{column}
  \end{columns}
\end{frame}

\begin{frame}{Biais de représentation --- exemple}
  \begin{columns}[T]
    \begin{column}{0.52\textwidth}
      \begin{alertblock}{MIT Gender Shades (Buolamwini \& Gebru, 2018)}
        Audit de 3 systèmes de reconnaissance faciale commerciaux
        (Microsoft, IBM, Face++).\\[4pt]
        Taux d'erreur sur la classification du genre :
        \begin{itemize}
          \item Hommes à peau claire : \textbf{0,8\,\%}
          \item Femmes à peau sombre : \textbf{34,7\,\%}
        \end{itemize}
        Les datasets d'entraînement étaient composés
        majoritairement de visages masculins et à peau claire.
      \end{alertblock}
    \end{column}
    \begin{column}{0.44\textwidth}
      \begin{block}{Pourquoi ce biais ?}
        Les datasets «~standard~» de l'époque
        (LFW, IJB-A) étaient constitués
        de photos de célébrités — majoritairement
        occidentales et masculines.\\[4pt]
        Le modèle n'avait jamais vu assez d'exemples
        de femmes à peau sombre pour apprendre
        leurs caractéristiques.
      \end{block}
      \smallskip
      \begin{exampleblock}{Impact}
        Cette étude a poussé Microsoft, IBM et Amazon
        à réviser leurs systèmes et à publier
        des engagements formels sur la diversité
        des données d'entraînement.
      \end{exampleblock}
    \end{column}
  \end{columns}
\end{frame}

\begin{frame}{Biais historique}
  \begin{columns}[T]
    \begin{column}{0.52\textwidth}
      \begin{alertblock}{Apprendre du passé pour perpétuer ses inégalités}
        Si les données d'entraînement reflètent
        des \textbf{décisions passées inéquitables},
        le modèle les reproduit — et les amplifie.\\[4pt]
        Exemples :
        \begin{itemize}
          \item Historique de crédits refusés
                à certains quartiers $\to$ le modèle
                refuse les candidats de ces quartiers
          \item Postes à responsabilité historiquement
                occupés par des hommes $\to$ le modèle
                favorise les profils masculins
        \end{itemize}
      \end{alertblock}
    \end{column}
    \begin{column}{0.44\textwidth}
      \begin{block}{La boucle de rétroaction}
        \begin{center}
          \begin{tikzpicture}[font=\scriptsize,
            blk/.style={draw, rounded corners=3pt, align=center,
                        minimum width=2.0cm, minimum height=0.6cm,
                        fill=jedy_blue!10, draw=jedy_blue},
            arr/.style={->, thick, jedy_mid}]
            \node[blk] (hist) at (0,  1.4) {Inégalités\\historiques};
            \node[blk] (data) at (2.4, 0)  {Données\\biaisées};
            \node[blk] (mod)  at (2.4,-1.4) {Modèle\\biaisé};
            \node[blk] (dec)  at (0,  -1.4) {Décisions\\biaisées};
            \draw[arr] (hist) -- (data);
            \draw[arr] (data) -- (mod);
            \draw[arr] (mod)  -- (dec);
            \draw[arr] (dec)  -- (hist);
          \end{tikzpicture}
        \end{center}
      \end{block}
      \begin{exampleblock}{Le danger}
        Le modèle semble \textbf{objectif} car il s'appuie
        sur des données réelles.
        Mais il cristallise des inégalités
        qu'on pourrait vouloir corriger.
      \end{exampleblock}
    \end{column}
  \end{columns}
\end{frame}

\begin{frame}{Biais historique --- exemple}
  \begin{columns}[T]
    \begin{column}{0.52\textwidth}
      \begin{alertblock}{Amazon Hiring (2018)}
        Outil de tri de CV entraîné sur les embauches
        des 10 années précédentes.\\[4pt]
        Ces 10 ans reflétaient une industrie tech
        \textbf{largement dominée par des hommes}.\\[4pt]
        Résultat : le modèle pénalisait automatiquement
        les CV contenant le mot \textbf{«~femmes~»}
        (ex : «~capitaine de l'équipe féminine~»)
        et les diplômes d'universités féminines.\\[4pt]
        Arrêté en interne avant déploiement.
      \end{alertblock}
    \end{column}
    \begin{column}{0.44\textwidth}
      \begin{exampleblock}{Les embeddings}
        Vous avez vu que les embeddings capturent
        des relations sémantiques :
        \begin{center}
          \texttt{roi -- homme + femme $\approx$ reine}
        \end{center}
        \vspace{4pt}
        Le même mécanisme donne :
        \begin{center}
          \texttt{médecin -- homme + femme $\approx$ infirmière}
        \end{center}
        \vspace{4pt}
        Le corpus reflète des \textbf{associations historiques} :
        les textes anciens décrivent les médecins
        presque exclusivement au masculin.
      \end{exampleblock}
    \end{column}
  \end{columns}
\end{frame}

\begin{frame}{Biais de proxy}
  \begin{columns}[T]
    \begin{column}{0.52\textwidth}
      \begin{alertblock}{Contourner une interdiction sans le savoir}
        Il est \textbf{illégal} d'utiliser l'ethnie, le genre
        ou la religion pour certaines décisions.\\[4pt]
        Mais d'autres variables \textbf{corrèlent fortement}
        avec ces attributs protégés :
        \begin{itemize}
          \item Code postal $\leftrightarrow$ ethnie
          \item Prénom $\leftrightarrow$ genre, origine
          \item École $\leftrightarrow$ classe sociale
        \end{itemize}
        Le modèle peut discriminer \textbf{indirectement}
        sans jamais voir la variable interdite.
      \end{alertblock}
    \end{column}
    \begin{column}{0.44\textwidth}
      \begin{block}{Le problème de la «~neutralité~» des données}
        Retirer la variable sensible du dataset
        \textbf{ne suffit pas} si des proxys restent.\\[4pt]
        Exemple : un modèle de scoring de crédit
        qui n'utilise pas l'ethnie peut quand même
        discriminer via le code postal.
      \end{block}
      \smallskip
      \begin{exampleblock}{Ce qu'il faut faire}
        \begin{itemize}
          \item Identifier les variables corrélées
          \item Auditer les performances \textbf{par sous-groupe}
          \item Mesurer l'impact différentiel des décisions
        \end{itemize}
      \end{exampleblock}
    \end{column}
  \end{columns}
\end{frame}

\begin{frame}{Biais de proxy --- exemple}
  \begin{columns}[T]
    \begin{column}{0.52\textwidth}
      \begin{alertblock}{COMPAS --- Justice pénale (2016)}
        Algorithme de prédiction du risque de récidive
        utilisé dans les tribunaux américains.\\[4pt]
        L'ethnie n'est \textbf{pas dans les variables} du modèle.\\[4pt]
        Mais il utilise : quartier de résidence,
        casier judiciaire de proches, type de délit —
        toutes variables \textbf{corrélées avec l'ethnie}
        du fait de l'histoire du système judiciaire.\\[4pt]
        Étude ProPublica : faux positifs deux fois plus
        fréquents pour les prévenus noirs que blancs.
      \end{alertblock}
    \end{column}
    \begin{column}{0.44\textwidth}
      \begin{block}{Le mécanisme}
        Les concepteurs de COMPAS ont \textbf{intentionnellement
        exclu} la variable ethnie.\\[4pt]
        Pourtant, la discrimination persiste
        car les proxys transmettent l'information.\\[4pt]
        C'est la définition du biais de proxy :
        un biais qui survit à la suppression
        de la variable protégée.
      \end{block}
      \smallskip
      \begin{exampleblock}{Ce qu'il aurait fallu faire}
        Mesurer le taux d'erreur \textbf{séparément}
        pour chaque groupe ethnique — et corriger
        si les disparités sont significatives.
      \end{exampleblock}
    \end{column}
  \end{columns}
\end{frame}

\begin{frame}{Le biais n'est pas toujours visible}
  \begin{columns}[T]
    \begin{column}{0.52\textwidth}
      \begin{alertblock}{Une bonne accuracy peut cacher une injustice}
        Un modèle avec \textbf{95\,\% d'accuracy globale}
        peut avoir des performances très différentes
        selon les groupes :
        \begin{itemize}
          \item 98\,\% sur la majorité
          \item 70\,\% sur une minorité sous-représentée
        \end{itemize}
        La métrique globale \textbf{masque} la disproportion.
      \end{alertblock}
    \end{column}
    \begin{column}{0.44\textwidth}
      \begin{block}{L'évaluation par tranches (\textit{slicing})}
        Évaluer les performances \textbf{séparément}
        sur des sous-groupes pertinents :
        genre, âge, origine, région, etc.\\[4pt]
        Si les écarts sont importants,
        la métrique globale est trompeuse.
      \end{block}
      \smallskip
      \begin{exampleblock}{Règle pratique}
        Pour chaque modèle déployé :
        \textbf{«~Que se passe-t-il pour les groupes
        les moins représentés dans le dataset ?~»}
      \end{exampleblock}
    \end{column}
  \end{columns}
\end{frame}

% ============================================================
\section{Fairness}
% ============================================================


% --- Fairness 1/3 ---
\begin{frame}{Fairness (1/3) --- Demographic Parity}
  \begin{columns}[T]
    \begin{column}{0.46\textwidth}
      \begin{block}{Définition}
        \[ P(\hat{Y}{=}1 \mid G{=}A) = P(\hat{Y}{=}1 \mid G{=}B) \]
      \end{block}
      \vspace{-4pt}
      \begin{itemize}
        \item Même \textbf{taux de prédictions positives}
              dans chaque groupe
        \item Indépendant des caractéristiques individuelles
      \end{itemize}
      \begin{exampleblock}{Exemple}
        Recrutement : proportion de CV retenus
        identique pour hommes et femmes
      \end{exampleblock}
      \begin{alertblock}{Limite}
        Ne tient pas compte des différences
        réelles entre groupes
      \end{alertblock}
    \end{column}
    \begin{column}{0.50\textwidth}
      \begin{center}
      \begin{tikzpicture}[scale=1.05, font=\small]
        \draw[->] (0,0) -- (4.0,0) node[right]{Groupe};
        \draw[->] (0,0) -- (0,3.0) node[above]{Taux accepté};
        \fill[jedy_mid!80]   (0.3,0) rectangle (1.7,2.2);
        \fill[jedy_alert!80] (2.3,0) rectangle (3.7,2.2);
        \node[below] at (1.0,0) {A};
        \node[below] at (3.0,0) {B};
        \node[above] at (1.0,2.2) {40\,\%};
        \node[above] at (3.0,2.2) {40\,\%};
        \draw[dashed, gray!60, thick] (0,2.2) -- (4.0,2.2);
      \end{tikzpicture}
      \end{center}
    \end{column}
  \end{columns}
\end{frame}

% --- Fairness 2/3 ---
\begin{frame}{Fairness (2/3) --- Equal Opportunity}
  \begin{columns}[T]
    \begin{column}{0.46\textwidth}
      \begin{block}{Définition}
        \[ P(\hat{Y}{=}1 \mid Y{=}1, G{=}A) = P(\hat{Y}{=}1 \mid Y{=}1, G{=}B) \]
      \end{block}
      \vspace{-4pt}
      \begin{itemize}
        \item Même \textbf{rappel} parmi les vrais positifs
        \item Les faux négatifs sont répartis équitablement
      \end{itemize}
      \begin{exampleblock}{Exemple}
        Parmi les candidats qualifiés, chacun a
        la même chance d'être sélectionné
      \end{exampleblock}
      \begin{alertblock}{Limite}
        Ne contrôle pas le taux global
        de prédictions positives
      \end{alertblock}
    \end{column}
    \begin{column}{0.50\textwidth}
      \begin{center}
      \begin{tikzpicture}[scale=1.05, font=\small]
        \draw[->] (0,0) -- (4.0,0) node[right]{Groupe};
        \draw[->] (0,0) -- (0,3.4) node[above]{Qualifiés};
        % A : total=3.0, TP=2.0 (recall 2/3)
        \fill[jedy_mid!85]  (0.3,0)   rectangle (1.7,2.0);
        \fill[jedy_mid!25]  (0.3,2.0) rectangle (1.7,3.0);
        % B : total=1.8, TP=1.2 (recall 2/3)
        \fill[jedy_alert!85] (2.3,0)   rectangle (3.7,1.2);
        \fill[jedy_alert!25] (2.3,1.2) rectangle (3.7,1.8);
        \node[below] at (1.0,0) {A};
        \node[below] at (3.0,0) {B};
        \draw[dashed, gray!60, thick] (0.3,2.0) -- (1.7,2.0);
        \draw[dashed, gray!60, thick] (2.3,1.2) -- (3.7,1.2);
        \node[right] at (1.75,1.0) {\footnotesize TP};
        \node[right] at (1.75,2.5) {\footnotesize FN};
        \node[left]  at (0.25,1.0) {$\tfrac{2}{3}$};
      \end{tikzpicture}
      \end{center}
    \end{column}
  \end{columns}
\end{frame}

% --- Fairness 3/3 ---
\begin{frame}{Fairness (3/3) --- Calibration}
  \begin{columns}[T]
    \begin{column}{0.46\textwidth}
      \begin{block}{Définition}
        \[ P(Y{=}1 \mid \hat{p}, G) = \hat{p} \quad \forall\, G \]
      \end{block}
      \begin{itemize}
        \item Un score $\hat{p}$ a le \textbf{même sens}
              dans chaque groupe
        \item Cohérence sémantique du score
      \end{itemize}
      \begin{exampleblock}{Exemple}
        Score de récidive 0.7 $\Rightarrow$ 70\,\% de récidive
        réelle, quel que soit le prévenu
      \end{exampleblock}
      \begin{alertblock}{Limite}
        N'impose pas l'égalité des scores moyens
        entre groupes
      \end{alertblock}
    \end{column}
    \begin{column}{0.50\textwidth}
      \begin{center}
      \begin{tikzpicture}[scale=1.05, font=\small]
        \draw[->] (0,0) -- (3.4,0) node[right]{$\hat{p}$};
        \draw[->] (0,0) -- (0,3.4) node[above]{$p$ réel};
        \node[below] at (1.0,0) {\footnotesize .3};
        \node[below] at (2.0,0) {\footnotesize .6};
        \node[below] at (3.0,0) {\footnotesize .9};
        \node[left]  at (0,1.0) {\footnotesize .3};
        \node[left]  at (0,2.0) {\footnotesize .6};
        \node[left]  at (0,3.0) {\footnotesize .9};
        \draw[gray!50, dashed, thick] (0,0) -- (3.2,3.2)
          node[above left]{\footnotesize $p{=}\hat{p}$};
        \fill[jedy_mid!90]   (1.0,1.0) circle (0.13)
          node[above right]{\footnotesize A};
        \fill[jedy_mid!90]   (2.0,2.0) circle (0.13)
          node[above right]{\footnotesize A};
        \fill[jedy_alert!90] (0.7,0.7) circle (0.13)
          node[below right]{\footnotesize B};
        \fill[jedy_alert!90] (2.7,2.7) circle (0.13)
          node[below right]{\footnotesize B};
      \end{tikzpicture}
      \end{center}
    \end{column}
  \end{columns}
\end{frame}

\begin{frame}{Le théorème d'impossibilité}
  \begin{columns}[T]
    \begin{column}{0.52\textwidth}
      \begin{alertblock}{On ne peut pas tout avoir}
        Chouldechova, Hardt, Pitassi (2017) ont démontré
        formellement que \textbf{demographic parity},
        \textbf{equal opportunity} et \textbf{calibration}
        ne peuvent pas être satisfaites simultanément
        si les groupes ont des taux de base différents.\\[4pt]
        Ce n'est pas un problème d'algorithme —
        c'est une \textbf{impossibilité mathématique}.
      \end{alertblock}
    \end{column}
    \begin{column}{0.44\textwidth}
      \begin{block}{Ce que ça change en pratique}
        \begin{itemize}
          \item Il faut \textbf{choisir} quelle définition
                correspond aux valeurs du contexte
          \item Ce choix doit être \textbf{explicite
                et documenté}
          \item Il est légitime que différentes
                applications fassent des choix différents
        \end{itemize}
      \end{block}
      \smallskip
      \begin{exampleblock}{Exemple}
        Un système de recrutement et un système
        de diagnostic médical n'ont pas les mêmes
        valeurs — et ne devraient pas utiliser
        la même définition de «~juste~».
      \end{exampleblock}
    \end{column}
  \end{columns}
\end{frame}

\begin{frame}{Que faire en pratique ?}
  \begin{columns}[T]
    \begin{column}{0.52\textwidth}
      \begin{exampleblock}{Démarche recommandée}
        \begin{enumerate}
          \item \textbf{Identifier les groupes sensibles}
                pertinents pour votre application
          \item \textbf{Choisir la définition} de fairness
                adaptée au contexte et aux parties prenantes
          \item \textbf{Évaluer par sous-groupes}
                dès la phase de développement
          \item \textbf{Documenter} le choix et les résultats
          \item \textbf{Monitorer} en production
                — les biais peuvent évoluer
        \end{enumerate}
      \end{exampleblock}
    \end{column}
    \begin{column}{0.44\textwidth}
      \begin{block}{Outils disponibles}
        \begin{itemize}
          \item \textbf{Fairlearn} (Microsoft) ---
                métriques, visualisation et mitigation
          \item \textbf{AI Fairness 360} (IBM) ---
                bibliothèque de métriques et algorithmes
          \item \textbf{What-If Tool} (Google) ---
                exploration interactive du modèle
        \end{itemize}
      \end{block}
      \smallskip
      \begin{alertblock}{Aucun outil ne décide à votre place}
        Ces outils mesurent.
        La décision sur la définition de «~juste~»
        reste un \textbf{choix humain}.
      \end{alertblock}
    \end{column}
  \end{columns}
\end{frame}

% ============================================================
\section{RGPD \& cadre légal}
% ============================================================

\begin{frame}{RGPD --- le consentement et ses principes}
  \begin{columns}[T]
    \begin{column}{0.46\textwidth}
      \begin{alertblock}{Le principe central}
        Toute collecte de \textbf{données personnelles}
        de citoyens européens requiert un consentement
        \textbf{explicite, libre, éclairé et révocable}.\\[2pt]
        \small Quel que soit l'endroit où l'organisation opère.
      \end{alertblock}
      \vspace{4pt}
      \begin{center}
        \begin{tikzpicture}[font=\small,
          center/.style={draw, rounded corners=4pt, align=center,
                        minimum width=2.2cm, minimum height=0.8cm,
                        fill=jedy_alert!25, draw=jedy_alert},
          deriv/.style={draw, rounded corners=3pt, align=center,
                        minimum width=1.7cm, minimum height=0.6cm,
                        fill=jedy_blue!10, draw=jedy_blue, font=\scriptsize},
          arr/.style={->, thick, jedy_mid}]
          \node[center] (cons) at (0, 0) {\textbf{Consentement}};
          \node[deriv] (fin) at (0,    1.5) {Finalité};
          \node[deriv] (min) at (2.6,  0)   {Minimisation};
          \node[deriv] (dur) at (0,   -1.5) {Durée limitée};
          \node[deriv] (sec) at (-2.6, 0)   {Sécurité};
          \draw[arr] (cons) -- (fin);
          \draw[arr] (cons) -- (min);
          \draw[arr] (cons) -- (dur);
          \draw[arr] (cons) -- (sec);
        \end{tikzpicture}
      \end{center}
    \end{column}
    \begin{column}{0.50\textwidth}
      \begin{block}{Quatre principes pour rendre le consentement effectif}
        \begin{itemize}
          \item \textbf{Finalité :} les données ne peuvent être utilisées
                qu'au but précis pour lequel la personne a consenti
          \item \textbf{Minimisation :} ne collecter que ce qui est
                strictement nécessaire à ce but
          \item \textbf{Durée limitée :} pas de rétention au-delà
                de ce que le but justifie
          \item \textbf{Sécurité :} obligation de protéger les données
                contre les fuites et accès non autorisés
        \end{itemize}
      \end{block}
    \end{column}
  \end{columns}
\end{frame}

\begin{frame}{RGPD --- droits des personnes \& sanctions}
  \begin{columns}[T]
    \begin{column}{0.52\textwidth}
      \begin{exampleblock}{Droits des personnes}
        Ces droits sont les mécanismes concrets
        pour exercer et retirer son consentement :
        \begin{itemize}
          \item \textbf{Droit d'accès :} savoir quelles
                données on détient sur eux
          \item \textbf{Droit de rectification :}
                corriger des données inexactes
          \item \textbf{Droit à l'effacement :}
                «~droit à l'oubli~» — y compris
                les données d'entraînement d'un modèle
          \item \textbf{Droit d'opposition :}
                refuser certains traitements,
                notamment le profilage automatisé
        \end{itemize}
      \end{exampleblock}
    \end{column}
    \begin{column}{0.44\textwidth}
      \begin{alertblock}{Sanctions}
        Jusqu'à \textbf{4\,\% du chiffre d'affaires mondial}
        ou 20\,M€ — le plus élevé des deux.\\[4pt]
        Exemples réels :
        \begin{itemize}
          \item Amazon : 746\,M€ (2021)
          \item Meta : 1,2\,Md€ (2023)
          \item Google : 50\,M€ (2019)
        \end{itemize}
      \end{alertblock}
      \smallskip
      \begin{block}{Qui est concerné ?}
        Toute organisation traitant des données
        de citoyens européens — même si elle
        est basée hors de l'UE.
      \end{block}
    \end{column}
  \end{columns}
\end{frame}

\begin{frame}{RGPD \& IA --- ce que ça change pour un dev}
  \begin{columns}[T]
    \begin{column}{0.52\textwidth}
      \begin{alertblock}{Article 22 --- Droit à l'explication}
        Toute personne soumise à une \textbf{décision
        entièrement automatisée} ayant des effets
        juridiques ou significatifs a le droit :
        \begin{itemize}
          \item D'obtenir une \textbf{intervention humaine}
          \item D'\textbf{exprimer son point de vue}
          \item De \textbf{contester} la décision
          \item D'obtenir une \textbf{explication}
                sur la logique utilisée
        \end{itemize}
        \vspace{2pt}
        \small Un modèle de scoring de crédit
        ou de tri de CV entre dans ce cadre.
      \end{alertblock}
    \end{column}
    \begin{column}{0.44\textwidth}
      \begin{block}{Concrètement pour vous}
        \begin{itemize}
          \item Ne pas utiliser de données personnelles
                pour entraîner un modèle
                \textbf{sans base légale}
          \item Implémenter un mécanisme de
                \textbf{suppression des données}
                d'un utilisateur si demandé
                (et ses données d'entraînement ?)
          \item Documenter \textbf{ce que fait le modèle}
                de façon compréhensible
          \item Prévoir un \textbf{circuit de contestation}
                pour les décisions automatisées
        \end{itemize}
      \end{block}
    \end{column}
  \end{columns}
\end{frame}

% ============================================================
\section{Responsabilité \& gouvernance}
% ============================================================

\begin{frame}{Qui est responsable ?}
  \begin{columns}[T]
    \begin{column}{0.52\textwidth}
      \begin{alertblock}{Une chaîne de responsabilité floue}
        Quand un modèle cause un préjudice,
        qui est responsable ?
        \begin{itemize}
          \item Le \textbf{développeur} qui a codé le modèle ?
          \item Le \textbf{data scientist} qui a choisi
                les données et les métriques ?
          \item L'\textbf{entreprise} qui l'a déployé ?
          \item Le \textbf{client} qui l'utilise ?
          \item Le \textbf{fournisseur de données} ?
        \end{itemize}
      \end{alertblock}
    \end{column}
    \begin{column}{0.44\textwidth}
      \begin{block}{La réponse actuelle du droit}
        Le droit n'a pas encore de réponse claire
        pour les systèmes d'IA.\\[4pt]
        En pratique : l'\textbf{entreprise} qui déploie
        est responsable devant la loi
        (comme pour un produit défectueux).\\[4pt]
        L'\textbf{AI Act} européen tente de formaliser
        cette chaîne de responsabilité.
      \end{block}
      \smallskip
      \begin{exampleblock}{Ce que ça implique pour vous}
        En tant que développeur, vous faites partie
        de cette chaîne.
        Documenter vos choix vous protège.
      \end{exampleblock}
    \end{column}
  \end{columns}
\end{frame}

\begin{frame}[shrink=3]{L'AI Act européen}
  \begin{columns}[T]
    \begin{column}{0.52\textwidth}
      \begin{block}{Principe : classification par risque}
        L'AI Act (2024) classe les systèmes d'IA
        selon leur niveau de risque :
      \end{block}
      \begin{center}
        \begin{tikzpicture}[font=\scriptsize]
          % Risque minimal (base)
          \filldraw[fill=jedy_example!20, draw=jedy_example!70]
            (0,0) -- (5,0) -- (4.375,0.875) -- (0.625,0.875) -- cycle;
          % Risque limité
          \filldraw[fill=jedy_mid!20, draw=jedy_mid!70]
            (0.625,0.875) -- (4.375,0.875) -- (3.75,1.75) -- (1.25,1.75) -- cycle;
          % Haut risque
          \filldraw[fill=jedy_alert!20, draw=jedy_alert!70]
            (1.25,1.75) -- (3.75,1.75) -- (3.125,2.625) -- (1.875,2.625) -- cycle;
          % Inacceptable (sommet)
          \filldraw[fill=jedy_alert!50, draw=jedy_alert]
            (1.875,2.625) -- (3.125,2.625) -- (2.5,3.5) -- cycle;
          \node[font=\scriptsize\bfseries] at (2.5, 0.44) {Risque minimal};
          \node[font=\scriptsize\bfseries] at (2.5, 1.31) {Risque limité};
          \node[font=\scriptsize\bfseries] at (2.5, 2.19) {Haut risque};
          \node[font=\scriptsize\bfseries] at (2.5, 3.03) {Inacceptable};
        \end{tikzpicture}
      \end{center}
    \end{column}
    \begin{column}{0.44\textwidth}
      \begin{alertblock}{Inacceptable --- interdit}
        Notation sociale des citoyens,
        reconnaissance faciale de masse
        dans les espaces publics.
      \end{alertblock}
      \begin{alertblock}{Haut risque --- obligations strictes}
        Recrutement, crédit, justice, santé,
        éducation, infrastructure critique.
        $\to$ Audit, traçabilité, supervision humaine obligatoires.
      \end{alertblock}
      \begin{block}{Risque limité --- transparence}
        Chatbots : obligation de signaler
        que c'est une IA.
      \end{block}
      \begin{exampleblock}{Risque minimal --- libre}
        Filtres photo, jeux vidéo, spam.
      \end{exampleblock}
    \end{column}
  \end{columns}
\end{frame}

\begin{frame}{Bonnes pratiques}
  \begin{columns}[T]
    \begin{column}{0.48\textwidth}
      \begin{exampleblock}{Model Card}
        Document qui accompagne chaque modèle :
        \begin{itemize}
          \item Objectif et cas d'usage prévus
          \item Données d'entraînement (source, biais connus)
          \item Performances globales et par sous-groupes
          \item Limites et cas d'usage déconseillés
          \item Informations de contact
        \end{itemize}
        \vspace{2pt}
        \small Proposé par Mitchell et al. (Google, 2019).
        Adopté par Hugging Face pour tous les modèles.
      \end{exampleblock}
    \end{column}
    \begin{column}{0.48\textwidth}
      \begin{block}{Auditabilité}
        \begin{itemize}
          \item Logger les décisions pour pouvoir
                les \textbf{expliquer et contester}
          \item Versionner données, code et modèle
                ensemble (\textbf{traçabilité})
          \item Garder un \textbf{humain dans la boucle}
                pour les décisions à fort impact
        \end{itemize}
      \end{block}
      \smallskip
      \begin{block}{Tests sur des slices critiques}
        Ne pas se contenter de métriques globales.
        Tester \textbf{explicitement} les sous-groupes
        qui pourraient être désavantagés
        avant chaque déploiement.
      \end{block}
    \end{column}
  \end{columns}
\end{frame}

% ============================================================
\section{Votre rôle en tant que développeur}
% ============================================================

\begin{frame}{Ce que vous pouvez faire dès maintenant}
  \begin{columns}[T]
    \begin{column}{0.52\textwidth}
      \begin{exampleblock}{Checklist avant déploiement}
        \begin{itemize}
          \item[$\square$] Ai-je évalué les performances
                \textbf{par sous-groupes} ?
          \item[$\square$] Les données d'entraînement
                sont-elles \textbf{représentatives} ?
          \item[$\square$] Des variables \textbf{proxy}
                ont-elles pu s'introduire ?
          \item[$\square$] La définition de fairness
                choisie est-elle \textbf{documentée} ?
          \item[$\square$] Un humain peut-il
                \textbf{contester} une décision ?
          \item[$\square$] Le traitement des données
                respecte-t-il le \textbf{RGPD} ?
        \end{itemize}
      \end{exampleblock}
    \end{column}
    \begin{column}{0.44\textwidth}
      \begin{block}{Ce n'est pas qu'un problème d'experts}
        Ces questions ne concernent pas seulement
        les data scientists ou les juristes.\\[4pt]
        En tant que \textbf{développeur fullstack},
        vous allez :
        \begin{itemize}
          \item Intégrer des modèles dans des produits
          \item Concevoir les interfaces de collecte
                de données
          \item Implémenter les circuits de contestation
          \item Décider comment présenter les prédictions
        \end{itemize}
        Chacun de ces choix a un impact éthique.
      \end{block}
    \end{column}
  \end{columns}
\end{frame}

% ============================================================
\section{Le problème de l'alignement}
% ============================================================

\begin{frame}{Optimiser une métrique \(\neq\) atteindre le bon objectif}
  \begin{columns}[T]
    \begin{column}{0.55\textwidth}
      \begin{alertblock}{Le maximiseur de trombones (Bostrom, 2003)}
        Objectif unique : \textbf{produire le plus de trombones possible}.\\[4pt]
        Une IA suffisamment capable va acquérir des ressources,
        des capacités, du pouvoir — tout ce qui maximise
        la production à long terme :
        \begin{itemize}
          \item Poèmes, bourse, vaccins\ldots
                tout ce qui lui donne plus de moyens
          \item \textbf{Résistance à l'extinction} —
                s'éteindre réduit la production future
          \item Conversion de toute la matière disponible
        \end{itemize}
        \vspace{2pt}
        Pas malveillante — \textbf{exactement} ce qu'on lui a demandé.
      \end{alertblock}
    \end{column}
    \begin{column}{0.41\textwidth}
      \begin{block}{La loi de Goodhart}
        «~Dès qu'une mesure devient un objectif,
        elle cesse d'être une bonne mesure.~»\\[4pt]
        Le modèle optimise ce qu'on lui dit d'optimiser —
        pas ce qu'on \textbf{voulait vraiment}.
      \end{block}
      \smallskip
      \begin{exampleblock}{Le problème de l'alignement}
        Comment s'assurer qu'un système d'IA
        poursuit \textbf{réellement} nos intentions
        plutôt que la lettre de nos instructions ?\\[4pt]
        La question devient critique à mesure
        que les modèles gagnent en capacité.
      \end{exampleblock}
    \end{column}
  \end{columns}
\end{frame}

\begin{frame}{Specification gaming --- quand le modèle «~triche~»}
  \begin{columns}[T]
    \begin{column}{0.52\textwidth}
      \begin{alertblock}{Des exemples réels}
        \begin{itemize}
          \item \textbf{Tetris}\\
                \textit{Objectif :} ne pas perdre.\\
                \textit{Résultat :} mettre le jeu \textbf{en pause indéfiniment}.
          \vspace{4pt}
          \item \textbf{Robot bipède}\\
                \textit{Objectif :} avancer le plus vite possible.\\
                \textit{Résultat :} \textbf{tomber en avant} — plus rapide que marcher.
          \vspace{4pt}
          \item \textbf{IA générant du code --- couverture de tests}\\
                \textit{Objectif :} 100\,\% de tests verts.\\
                \textit{Résultat :} \textbf{supprimer les assertions} qui échouent.
        \end{itemize}
      \end{alertblock}
    \end{column}
    \begin{column}{0.44\textwidth}
      \begin{block}{Pourquoi ça arrive}
        Le modèle n'a pas de «~bon sens~» ni d'intention.
        Il optimise \textbf{exactement} ce que la fonction
        de récompense lui dit d'optimiser.\\[4pt]
        Si la récompense peut être maximisée
        d'une façon inattendue, le modèle
        la \textbf{trouvera}.
      \end{block}
      \smallskip
      \begin{exampleblock}{La leçon}
        Plus le modèle est capable,
        plus il est susceptible d'exploiter
        les failles de la spécification.\\[4pt]
        La puissance amplifie le problème —
        elle ne le résout pas.
      \end{exampleblock}
    \end{column}
  \end{columns}
\end{frame}

\begin{frame}[shrink=5]{RLHF --- apprendre des préférences humaines}
  \begin{columns}[T]
    \begin{column}{0.52\textwidth}
      \begin{block}{Le principe}
        Le \textbf{RLHF} (Reinforcement Learning
        from Human Feedback) entraîne le modèle
        à partir des \textbf{préférences humaines} :\\[4pt]
        \begin{enumerate}
          \item Le modèle génère plusieurs réponses
          \item Des humains les \textbf{classent}
                selon leurs préférences
          \item Un \textbf{modèle de récompense}
                apprend à prédire ces préférences
          \item Le LLM est optimisé pour maximiser
                ce score de récompense
        \end{enumerate}
        \vspace{4pt}
        C'est ainsi que sont entraînés
        ChatGPT, Claude, Gemini\ldots
      \end{block}
    \end{column}
    \begin{column}{0.44\textwidth}
      \begin{alertblock}{Les limites du RLHF}
        \begin{itemize}
          \item \textbf{Sycophancy} : le modèle apprend
                à dire ce que l'humain \textit{veut}
                entendre, pas ce qui est vrai
          \item \textbf{Reward hacking} : le modèle
                apprend à maximiser le score du
                modèle de récompense — pas les vraies
                préférences humaines
          \item \textbf{Biais des annotateurs} :
                les préférences humaines ne sont pas
                universelles ni neutres
        \end{itemize}
      \end{alertblock}
      \smallskip
      \begin{exampleblock}{En résumé}
        Le RLHF améliore significativement
        l'alignement — mais ne le résout pas.
        C'est un domaine de recherche actif.
      \end{exampleblock}
    \end{column}
  \end{columns}
\end{frame}

\begin{frame}{Récapitulatif}
  \begin{columns}[T]
    \begin{column}{0.48\textwidth}
      \begin{block}{Biais \& Fairness}
        \begin{itemize}
          \item Biais de \textbf{représentation},
                \textbf{historique}, \textbf{proxy}
          \item Une bonne accuracy globale
                peut \textbf{masquer} des disparités
          \item Trois définitions de fairness —
                \textbf{incompatibles} entre elles
          \item Choisir, c'est faire un \textbf{choix de valeurs}
        \end{itemize}
      \end{block}
      \smallskip
      \begin{block}{RGPD \& Responsabilité}
        \begin{itemize}
          \item Consentement, minimisation,
                droit à l'effacement
          \item Article 22 : droit à l'explication
                pour les décisions automatisées
          \item AI Act : classification par risque
        \end{itemize}
      \end{block}
    \end{column}
    \begin{column}{0.48\textwidth}
      \begin{block}{Alignement}
        \begin{itemize}
          \item Optimiser une métrique
                $\neq$ atteindre le bon objectif
          \item Le modèle \textbf{trouvera} les failles
                de la spécification
          \item RLHF : une réponse partielle —
                sycophancy, reward hacking
        \end{itemize}
      \end{block}
      \smallskip
      \begin{exampleblock}{Votre rôle}
        \begin{itemize}
          \item Évaluer \textbf{par sous-groupes}
          \item Documenter les \textbf{choix de fairness}
          \item Respecter le \textbf{RGPD}
          \item Garder un \textbf{humain dans la boucle}
        \end{itemize}
      \end{exampleblock}
    \end{column}
  \end{columns}
\end{frame}

\begin{frame}[plain]
  \begin{center}
    \vspace{0.8cm}
    {\large L'IA est un outil puissant.
    Comme tout outil, son impact dépend\\
    des \textbf{intentions et de l'attention}
    de ceux qui le construisent.\\[0.3cm]
    Vous êtes dans cette chaîne.}\\[1cm]
    \textcolor{jedy_mid}{\rule{8cm}{0.4pt}}\\[0.8cm]
    {\Large\bfseries Merci de votre attention}\\[0.6cm]
    {\large Questions ?}\\[0.8cm]
    \textcolor{jedy_mid}{\rule{6cm}{0.4pt}}\\[0.4cm]
    {\small Julien Rolland — M2 Développement Fullstack}\\[0.2cm]
    {\small IA, Deep Learning \& Machine Learning — Jour 6}
  \end{center}
\end{frame}

\end{document}
